\section{Failure mode I: nonnegative residuals lose acceleration}
\label{sec:signed-star-monotone-failure}

Signed SOR succeeds because its over-relaxed step creates a controlled
opposite-sign reflection.  We now show, on exactly the same PPR instance and
under exactly the same semantic error, that preserving residual nonnegativity
forces the classical $1/\alpha$ scale.

\paragraph{Nonnegative-residual one-hop class $\mathcal M_+$.}
A member maintains the same exact invariant
$r=\gamma_\alpha e_o-Hp$, starts from $p=0$, and uses coordinate operations
\eqref{eq:signed-star-coordinate-primitive}.  Step sizes and coordinate order
may be adaptive, signed, or randomized, but every intermediate residual must
satisfy
\begin{equation}
  r\ge0.
  \label{eq:signed-star-nonnegative-residual}
\end{equation}
The output is $p$ and every operation is charged by the updated coordinate's
degree.  This class abstracts the order structure used by conservative local
push and monotone coordinate schemes.  It does not include SOR with
$\omega>1$ or FISTA extrapolation.

\subsection{A multiplicative center-deficit barrier}

Let
\begin{equation}
  q:=\pi_o-p_o=(H^{-1}r)_o
  \label{eq:signed-star-center-deficit}
\end{equation}
be the remaining center deficit.  Since $H^{-1}\ge0$ and $r\ge0$, we have
$q\ge0$.  On the star,
\begin{equation}
  (H^{-1})_{oo}=\frac{1}{1-c_\alpha^2},
  \qquad
  \kappa_\alpha:=1-c_\alpha^2
     =\frac{4\alpha}{(1+\alpha)^2}.
  \label{eq:signed-star-kappa}
\end{equation}
Dropping every noncenter contribution in $(H^{-1}r)_o$ gives
\begin{equation}
  q\ge\frac{r_o}{\kappa_\alpha},
  \qquad\text{hence}\qquad
  r_o\le\kappa_\alpha q.
  \label{eq:signed-star-center-residual-deficit}
\end{equation}

Consider a positive center update $p_o^+=p_o+\delta$ with $\delta>0$.
Because $H_{oo}=1$, its new center residual is $r_o^+=r_o-\delta$.
Nonnegativity forces $\delta\le r_o$, and therefore
\begin{equation}
  q^+=q-\delta
  \ge(1-\kappa_\alpha)q.
  \label{eq:signed-star-deficit-contraction-limit}
\end{equation}
A leaf operation does not change $p_o$ and hence does not change $q$; a
negative center operation only increases $q$.  Thus after $N_o^+$ positive
center operations, regardless of all intervening work,
\begin{equation}
  q_T\ge(1-\kappa_\alpha)^{N_o^+}\pi_o.
  \label{eq:signed-star-deficit-product}
\end{equation}

\begin{theorem}[Nonnegative-residual lower bound]
\label{thm:signed-star-monotone-lower}
Let $0<\alpha\le1/16$ and choose $m$ by
\eqref{eq:signed-star-size}.  Every successful path of every
$\mathcal M_+$ member satisfies
\begin{equation}
  N_o^+
  \ge
  \frac{\log\bigl(\pi_o/(\varepsilon_{\rm ppr}m)\bigr)}
       {-\log(1-\kappa_\alpha)}
  \ge\frac{3\log4}{16\alpha},
  \label{eq:signed-star-monotone-center-count}
\end{equation}
and hence
\begin{equation}
  W
  \ge mN_o^+
  \ge\frac{3\log4}
           {256\alpha\varepsilon_{\rm ppr}}
  =\Omega\!\left(
      \frac{1}{\alpha\varepsilon_{\rm ppr}}
    \right).
  \label{eq:signed-star-monotone-work}
\end{equation}
\end{theorem}

\begin{proof}
Semantic accuracy and $q_T\ge0$ imply
$q_T\le\varepsilon_{\rm ppr}m$.  Combine this with
\eqref{eq:signed-star-deficit-product}.  Since
$\pi_o\ge1/2$ and $\varepsilon_{\rm ppr}m\le1/8$, the logarithm in the
numerator is at least $\log4$.

For $\alpha\le1/16$,
$\kappa_\alpha\le64/289<1/4$.  Using
$-\log(1-x)\le x/(1-x)$ for $x\in(0,1)$ and
$\kappa_\alpha\le4\alpha$ gives
\begin{equation*}
  -\log(1-\kappa_\alpha)
  \le\frac43\kappa_\alpha
  \le\frac{16}{3}\alpha.
\end{equation*}
This proves the count bound.  Each positive center operation costs
$d_o=m\ge1/(16\varepsilon_{\rm ppr})$.
\end{proof}

The conclusion is a clean same-instance separation:
\begin{equation}
  \mathcal M_+:
  \Omega\!\left(\frac1{\alpha\varepsilon_{\rm ppr}}\right),
  \qquad
  \mathcal R_\pm:
  \Theta\!\left(\frac1{\sqrt\alpha\varepsilon_{\rm ppr}}\right).
  \label{eq:signed-star-sign-separation}
\end{equation}
These two displayed bounds concern the separately defined classes
$\mathcal M_+$ and $\mathcal R_\pm$; neither class contains the other.
Indeed, $\mathcal M_+$ permits arbitrary signed step sizes subject to the
residual constraint, whereas $\mathcal R_\pm$ requires
$\delta=\omega r_u$ with $0<\omega<2$ but permits signed residuals.
There is nevertheless a genuinely nested corollary.  Let
\begin{equation*}
 \mathcal R_\pm^{\rm nn}
 :=\{\text{runs in $\mathcal R_\pm$ with $r^t\geq0$ at every step}\}.
\end{equation*}
Then $\mathcal R_\pm^{\rm nn}\subset\mathcal M_+$, so
\Cref{thm:signed-star-monotone-lower} gives the
$\Omega(1/(\alpha\varepsilon_{\rm ppr}))$ lower bound for this subclass,
while optimal over-relaxed SOR in $\mathcal R_\pm$ attains the signed
$\Theta(1/(\sqrt\alpha\varepsilon_{\rm ppr}))$ scale.  In that nested
comparison, the decisive extra permission is a controlled signed reflection.

The theorem should not be broadened into an information-theoretic lower
bound for local PPR.  Exact block elimination, response methods, or an output
postprocessor with a charged stronger primitive are outside
$\mathcal M_+$.  Even within coordinate algorithms, the theorem assumes the
exact nonnegative residual invariant and output $p$; it does not classify
every conceivable nonlinear sparse representation.
